
\section{Tổng quan phương pháp}

\subsection{Bài toán và động lực nghiên cứu}

Bài toán phân bổ và chuyển ngang hàng hóa đa nhà máy (Multi-Facility Allocation and Proactive Lateral Transshipment -- MFA-PLT) được mô tả bằng mô hình tối ưu hóa có cấu trúc hỗn hợp nguyên--thực (MILP). Mô hình bao gồm:

\begin{itemize}
  \item \textbf{Biến quyết định liên tục và nguyên:} lượng phân bổ OA ($\QOA_{p,i,t}$), số case đầy đủ ($\qOA_{p,i,t}$), phần dư ($\rOA_{p,i,t}$); tương tự cho PLT.
  \item \textbf{Ràng buộc phi tuyến tính tiềm ẩn:} hàm $\pos{\cdot}$ (phần dương) xuất hiện trong hàm mục tiêu.
  \item \textbf{Không gian nghiệm lớn:} số biến tăng theo $|P| \times |L|^2 \times |T|$.
  \item \textbf{Tính phụ thuộc thời gian:} tồn kho $I_{p,i,t}$ phụ thuộc vào quyết định trong quá khứ thông qua lead time $LT_i^{OA}$ và $LT_{i,j}^{PLT}$.
\end{itemize}

Những đặc điểm này khiến cho việc giải chính xác bằng MILP solver trở nên kém khả thi khi quy mô bài toán tăng lên. Do đó, một phương pháp meta-heuristic kết hợp được đề xuất.

\subsection{Kiến trúc tổng thể}

Phương pháp đề xuất kết hợp hai meta-heuristic:

\begin{enumerate}
  \item \textbf{Genetic Algorithm (GA):} đóng vai trò khám phá không gian nghiệm rộng (global exploration), duy trì sự đa dạng quần thể và tránh rơi vào cực tiểu địa phương.
  \item \textbf{Adaptive Large Neighborhood Search (ALNS):} đóng vai trò khai thác cục bộ (local exploitation), tinh chỉnh từng nghiệm cá thể thông qua các toán tử phá hủy--tái tạo có trọng số thích nghi.
\end{enumerate}

\medskip

\noindent\textbf{Cơ chế tích hợp:} ALNS được nhúng vào GA như một \emph{bước local search} ngay sau bước giải mã (decode). Cụ thể, sau khi mỗi nhiễm sắc thể được giải mã thành nghiệm thực, ALNS được kích hoạt để cải thiện nghiệm đó trước khi đánh giá fitness. Cơ chế này còn được gọi là \emph{Lamarckian evolution} vì kinh nghiệm học được trong local search được phản ánh ngược lại vào nhiễm sắc thể.

\subsection{Phân rã bài toán theo sản phẩm}

Vì ràng buộc năng lực $\sum_{i \in L} \QOA_{p,i,t} = CAP_{p,t}$ tách rời theo từng sản phẩm $p$, bài toán có thể được phân rã thành $|P|$ bài toán con độc lập, mỗi bài toán con tối ưu hóa phân bổ cho một sản phẩm. GA-ALNS được áp dụng song song hoặc tuần tự cho từng sản phẩm.

\begin{figure}[H]
\centering
\begin{tikzpicture}[
  node distance = 1.6cm,
  every node/.style = {font=\small}
]
  \node (p1) [process, minimum width=5cm] {Sản phẩm $p=1$: GA-ALNS};
  \node (p2) [process, below of=p1, minimum width=5cm] {Sản phẩm $p=2$: GA-ALNS};
  \node (dots) [below of=p2, yshift=0.3cm] {$\vdots$};
  \node (pp) [process, below of=dots, yshift=0.3cm, minimum width=5cm] {Sản phẩm $p=|P|$: GA-ALNS};
  \node (agg) [io, below of=pp, minimum width=5cm] {Tổng hợp $Z = \sum_p Z_p$};

  \draw[arrow] (p1) -- (p2);
  \draw[arrow] (p2) -- (dots);
  \draw[arrow] (dots) -- (pp);
  \draw[arrow] (pp) -- (agg);
\end{tikzpicture}
\caption{Phân rã bài toán theo từng sản phẩm}
\end{figure}

\subsection{Biểu diễn nhiễm sắc thể}

\subsubsection{Phiên bản biểu diễn}

Đối với mỗi sản phẩm $p$ đã được cố định, một nhiễm sắc thể $\mathcal{C}$ được định nghĩa như một ma trận:

\begin{equation}
  \mathcal{C} = \left(\QOA_{i,t}\right)_{i \in L,\, t \in T}
  \in \mathbb{Z}_+^{|L| \times |T|}
\end{equation}

trong đó $\QOA_{i,t}$ là \emph{tổng lượng} hàng được phân bổ từ nguồn trung tâm đến nhà máy $i$ tại giai đoạn $t$ (bỏ chỉ số $p$ vì đã cố định). Mỗi ô trong ma trận là một số nguyên không âm.

\subsubsection{Ràng buộc trên nhiễm sắc thể}

\begin{equation}
  \sum_{i \in L} \QOA_{i,t} = CAP_t, \quad \forall t \in T
  \tag{Ràng buộc năng lực -- cứng}
\end{equation}

\begin{equation}
  \QOA_{i,t} \geq 0, \quad \forall i \in L, t \in T
\end{equation}

\subsubsection{Biến suy ra (không encode)}

Các biến sau được suy ra trong quá trình giải mã, không phải encode trực tiếp:

\begin{align}
  \qOA_{i,t} &= \left\lfloor \frac{\QOA_{i,t}}{CP_{i,t}} \right\rfloor
  \tag{số case đầy đủ}\\[4pt]
  \rOA_{i,t} &= \QOA_{i,t} - CP_{i,t} \cdot \qOA_{i,t}
  \tag{phần dư OA}\\[4pt]
  \QPLT_{i,j,t} &= f(\QOA, I)
  \tag{suy ra từ trạng thái tồn kho}\\[4pt]
  I_{i,t} &= g(\QOA, \QPLT, BI_i, \Delta I)
  \tag{cân bằng tồn kho}
\end{align}

\subsubsection{Ví dụ minh họa -- nhiễm sắc thể với 2 nhà máy, 3 giai đoạn}

Giả sử $CAP_t = 100$ đơn vị cho mọi $t$. Một nhiễm sắc thể hợp lệ:

\[
\mathcal{C} =
\begin{pmatrix}
  \QOA_{1,1} & \QOA_{1,2} & \QOA_{1,3} \\
  \QOA_{2,1} & \QOA_{2,2} & \QOA_{2,3}
\end{pmatrix}
=
\begin{pmatrix}
  40 & 55 & 30 \\
  60 & 45 & 70
\end{pmatrix}
\]

Kiểm tra: $40+60=100$, $55+45=100$, $30+70=100$ $\checkmark$

\subsection{Hàm giải mã (Decode)}

\label{subsec:decode}

Hàm giải mã chuyển đổi nhiễm sắc thể $\mathcal{C}$ thành một nghiệm hoàn chỉnh gồm $(\QOA, \QPLT, I)$:

\begin{algorithm}[H]
\caption{Giải mã nhiễm sắc thể (Decode)}
\begin{algorithmic}[1]
\Require Nhiễm sắc thể $\mathcal{C} = (\QOA_{i,t})$, tham số mô hình
\Ensure Nghiệm $(\QOA, \QPLT, I)$

\State \textbf{Bước 1: Tính phần dư OA}
\For{$i \in L$, $t \in T$}
  \State $\qOA_{i,t} \gets \lfloor \QOA_{i,t} / CP_{i,t} \rfloor$
  \State $\rOA_{i,t} \gets \QOA_{i,t} - CP_{i,t} \cdot \qOA_{i,t}$
\EndFor

\State \textbf{Bước 2: Mô phỏng tồn kho và xác định PLT}
\State $I_{i,0} \gets BI_i$ \quad cho mọi $i \in L$
\For{$t = 1$ đến $|T|$}
  \State \textit{-- Tính tồn kho tạm thời trước PLT --}
  \State $I^{temp}_{i,t} \gets I_{i,t-1} + \QOA_{i,t - LT_i^{OA}} + \Delta I_{i,t}$
  \quad (nếu $t - LT_i^{OA} < 1$ thì $\QOA_{i,\cdot} = 0$)
  \State \textit{-- Xác định PLT theo heuristic ưu tiên thiếu hụt --}
  \If{$t \in T_j^{PLT}$ với $T_j^{PLT} = \{t' : t' \leq LT_j^{OA} - 1\}$}
    \State Xếp hạng nhà máy theo nhu cầu PLT: $\text{need}_i = \pos{L_{i,t} - I^{temp}_{i,t}}$
    \State Xếp hạng nhà máy theo khả năng cho: $\text{supply}_i = \pos{I^{temp}_{i,t} - L_{i,t}}$
    \State Phân bổ PLT từ nhà máy thừa sang nhà máy thiếu (Thuật toán~\ref{alg:plt_heuristic})
  \EndIf
  \State \textit{-- Cập nhật tồn kho sau PLT --}
  \State $I_{i,t} \gets I^{temp}_{i,t} + \sum_{j \neq i} \QPLT_{j,i,t-LT_{j,i}^{PLT}} - \sum_{j \neq i} \QPLT_{i,j,t}$
\EndFor

\State \Return $(\QOA, \QPLT, I)$
\end{algorithmic}
\end{algorithm}

\begin{algorithm}[H]
\caption{Heuristic phân bổ PLT (trong bước giải mã)}
\label{alg:plt_heuristic}
\begin{algorithmic}[1]
\Require $I^{temp}_{i,t}$, $L_{i,t}$, $CP_{j,t}$ (case-pack của nhà máy nhận)
\Ensure $\QPLT_{i,j,t}$ cho mọi cặp $(i,j)$

\State Xác định tập nhà máy thiếu: $S^- = \{i : I^{temp}_{i,t} < L_{i,t}\}$
\State Xác định tập nhà máy dư: $S^+ = \{i : I^{temp}_{i,t} > L_{i,t}\}$
\State Xếp hạng $S^-$ theo mức thiếu hụt giảm dần: $\delta_i^- = L_{i,t} - I^{temp}_{i,t}$
\State Xếp hạng $S^+$ theo mức dư thừa giảm dần: $\delta_j^+ = I^{temp}_{j,t} - L_{j,t}$

\For{$j \in S^+$ (theo thứ tự)}
  \State $avail_j \gets \delta_j^+$
  \For{$i \in S^-$ (theo thứ tự)}
    \State $\text{need}_i \gets \delta_i^-$
    \State $\text{transfer} \gets \min(avail_j, \text{need}_i)$
    \State $k \gets \lfloor \text{transfer} / CP_{i,t} \rfloor$ \quad \textit{-- làm tròn theo case-pack nhà máy nhận --}
    \State $\QPLT_{j,i,t} \gets k \cdot CP_{i,t}$
    \State $\qPLT_{j,i,t} \gets k$; \quad $\rPLT_{j,i,t} \gets 0$ \quad \textit{-- nếu chỉ cho phép case đầy đủ --}
    \State $avail_j \gets avail_j - \QPLT_{j,i,t}$
    \State $\delta_i^- \gets \delta_i^- - \QPLT_{j,i,t}$
  \EndFor
\EndFor

\State \Return $\QPLT$
\end{algorithmic}
\end{algorithm}

\subsection{Hàm tính Fitness}

Fitness của nhiễm sắc thể được định nghĩa là giá trị hàm mục tiêu $Z$ sau khi giải mã:

\begin{align}
  f(\mathcal{C}) &= Z(\QOA, \QPLT, I) \notag\\
  &= \underbrace{\sum_{i,t} \left( C^o_{i,t} \pos{I_{i,t}-U_{i,t}} + C^s_{i,t} \pos{L_{i,t}-I_{i,t}} + C^b_{i,t} \pos{-I_{i,t}} \right)}_{\text{chi phí tồn kho}} \notag\\
  &\quad + \underbrace{\sum_{i,t} C^p_{i,t} \cdot \ind{\rOA_{i,t}>0}}_{\text{phạt OA không đủ case}} \notag\\
  &\quad + \underbrace{\sum_{i,j \neq i,\, t \in T_j^{PLT}} \left( C^p_{i,j,t} \cdot \ind{\rPLT_{i,j,t}>0} + d_{i,j} \cdot TC \cdot \ind{\QPLT_{i,j,t}>0} \right)}_{\text{chi phí PLT}}
  \label{eq:fitness}
\end{align}

Vì mục tiêu là \textbf{tối thiểu hóa}, fitness tốt hơn tương ứng với giá trị $f$ nhỏ hơn.
